 \documentclass[11pt,fleqn]{report}

\usepackage{cml}
\usepackage{amssymb}
\usepackage{graphicx}
\usepackage{subcaption}
\usepackage{longtable}

\newcommand{\ModelName}{Doppler Radar Wind Pairs (DRWP)}
\newcommand{\ModelAuthor}{Brian Birmingham, Gary Turner}
\newcommand{\ModelAffiliation}{Odyssey Space Research}
\newcommand{\ModelDate}{April 2023}

\setcounter{tocdepth}{3}
\setcounter{secnumdepth}{3}

\begin{document}
\pagestyle{empty}
\titlePage
\pagenumbering{roman}
\tableofcontents % Print table of contents
\vfill

{\Large \textbf {Revision History}}

\begin{tabular}{| m{0.1\textwidth}| m{0.15\textwidth} | m{0.2\textwidth} | m{0.4\textwidth} |} \hline
 \textbf{Version} & \textbf{Date} & \textbf{Author} & \textbf{Purpose} \\
\hline
1 & April 2023 & Brian Birmingham,&  \\
  &            & Gary Turner      & Initial Version \\ \hline
\end{tabular}


\chapter{Introduction} % Make a "chapter" heading
\pagestyle{plain}
\pagenumbering{arabic} % Start text with arabic 1
The Doppler Radar Wind Pairs (DRWP) model uses a one dimensional interpolation
algorithm to obtain atmospheric values s a function of altitude; altitude is
used as the independent variable for all interpolations.
The values directly interpolated from the data file are:
\begin{itemize}
\item 2 or 3 components of the Cartesian wind velocity vector expressed in
the topocentric East-North-Up (ENU) frame; the vertical component (U) is
optional, the two horizontal components (E,N) are required.
\item atmospheric mass-density
\item atmospheric kinetic-temperature
\item atmospheric static-pressure
\end{itemize}

In addition, the model computes values based on these interpolated values,
and outputs:
\begin{itemize}
\item A 3-component Cartesian vector representing the wind velocity
expressed in the topocentric North-East-Down(NED) frame.
\item The azimuth angle representing the direction from which the wind is
blowing in the horizontal plane.
\item The azimuth angle representing the direction to which the wind is
blowing in the horizontal plane.
\item The magnitude of the wind velocity vector.
\item The speed of sound.
\end{itemize}

Additionally, the model provides optional capability to compute the average
wind velocity across a range of altitudes.

The DRWP model is part of the Common Model Library (cml).

\chapter{Requirements}
\begin{enumerate}
\item The model shall provide a simple mechanism for interpolating wind
  and atmospheric values from a provided data set. The model shall provide
  the capability to interpolate the data presented according to the
  following assumptions:
   \begin{itemize}
  \item The file will be binary with a pre-defined data format.
    \footnote{The model assumes a specific data layout within the binary
              data file in response to this requirement. The organization
              of data within the data file
              is outlined in section \ref{sec:input_data}
              and discussed in more detail in section
              \ref{sec:binary_read_struc}.}
  \item Altitude will be the independent variable for interpolations.
  \item The file will include calibrated points at a selection of altitudes
    for the following variables:
    \begin{itemize}
    \item East component of wind velocity.
    \item North (Topocentric-North) component of wind velocity.
    \item (optional) Vertical (Topocentric-Up) component of wind velocity.
    \item mass-density.
    \item kinetic temperature.
    \item pressure.
    \end{itemize}
  \item For altitudes outside the domain of the provided data, the model
    shall return the data associated with the closest calibrated data set.
  \item Altitude data will be monotonically increasing.
   \end{itemize}
\item The model shall output the interpolated values for the following
variables:
  \begin{enumerate}
  \item the 2 or 3 components of the Cartesian wind velocity as provided in
    the data file; these are expressed in the Topocentric ENU frame.
  \item A 3-component Cartesian vector representing the wind velocity
    expressed in the Topocentric NED frame.
  \item The azimuth angle representing the direction from which the wind is
    blowing in the horizontal plane.
  \item The azimuth angle representing the direction to which the wind is
    blowing in the horizontal plane.
  \item The magnitude of the wind velocity vector.
  \item The atmospheric mass-density
  \item The atmospheric kinetic-temperature
  \item The atmospheric static-pressure.
  \item The local speed of sound.
  \end{enumerate}
\item The model shall provide an On/Off mechanism that will prevent the model
  from being executed if it is not needed by the simulation.
  The model should run only if it is explicitly activated.
\end{enumerate}



\newpage
\chapter{Model Specification}
\section{Inheritance}\label{sec:Inheritance}

The DRWP model inherits directly from the \textit{SubscriptionBase} class
defined in the CML
subscription model (\textit{utilities/subscriptions}).
This allows for the model to be be selectively activated, and disabled when
not being used by the simulation.

The DRWP model defines an interpolation-table class
(\textit{DRWPTableLookup}) that inherits from the
\textit{SimpleTableLookup} class defined in the CML table-interpolation
model (\textit{utilites/table\_interp\_cpp}). This provides the
interpolation system, as well as extensive tools for checking data
consistency.


\section{Dependencies}
In addition to the dependencies associated with inheritance, this model also
requires the CML Message Handling model (\textit{utilities/cml\_message}).

\section{Related Models} \label{sec:related_models}
This model is supported by the CML Atmosphere-executive model
(\textit{environment/atmos/atmos\_exec}) which handles switching between
atmospheres and the associated initialization, activation, execution, and
deactivation processes.

\section{Classes}
The model has two classes:
\begin{itemize}
\item \textbf{DRWPTableLookup}
\item \textbf{LookupAtmosWinds}
\end{itemize}

\subsection{DRWPTableLookup}
This class is an extension of the \textit{SimpleTableLookup} class, which
already provides all necessary interpolation algorithms. This class adds
four data members and one method to enhance management of multiple tables by
the \textit{LookupAtmosWinds} class.

The data contained within an instance of \textit{DRWPTableLookup} represents
a single profile, with one value for each of the interpolated variables for
each of one set of calibrated altitude values.

An important consideration associated with this class is the visibility of
its members and methods. Everything added to \textit{SimpleTableLookup}
by this class is public, but the
instances of this class within this model are private to the
\textit{LookupAtmosWinds} class so their status as public does not
imply a direct connection to the user-interface.
From the user-interface perspective, they should be considered private.
Of course, additional instances of this class could be created, and the
members of those instances would be accessible to the user-interface, but
these would not be accessible by the \textit{LookupAtmosWinds} class so
doing so would be futile within the intenided usage of this model.

\subsubsection{DRWPTableLookup Members}
\begin{longtable}{|p{2.0cm}|p{3.4cm}|p{5cm}|p{1.0cm}|p{2.0cm}|}
\caption{Table Management Data}\\
\hline
\textbf{Type} & \textbf{Name} & \textbf{Purpose} & \textbf{Units} & \textbf{Protection}\\
\hline
std::string & drwpFileName & The name of the data file from which this table
was populated & none & public \\
\hline
unsigned int & profile\_number & The profile number used to identify this
profile within the data file from which this table was populated. & none &
public \\
\hline
bool & initialized\_with\_\hspace{0.5cm}vertical\_component & This data table includes
vertical wind data & none & public \\
\hline
bool & verified & This table has been checked for internal consistency &
none & public \\
\hline
\end{longtable}

\subsubsection{DRWPTableLookup::verify(...)}
This method is used to perform internal consistency checking of the
populated data. This is primarily used at model initialization to support
checking all tables before the run starts executing; if an error is found,
the execution is terminated so it can be resolved before the run gets
underway.

As with the data members, this method is public but it has no user-interface
because the class instances within the model are private. The intent is that
the \textit{LookupAtmosWinds} class calls this method.

If the verification passes, the member flag \textit{verified} is set to
true. This flag is queried by the \textit{LookupAtmosWinds} model when it
switches to using a new data table mid-run. If the flag is true, the table
initialization and verification can be bypassed. If the flag is still false,
the initialization and verification of the table will be executed at the
time the table is first used, which can lead to simulation termination if an
anomaly is found.

The method takes a single optional argument; this is a boolean value that
defaults to false if not provided. This argument is an instruction to leave
the table active after verification is complete.


\subsection{LookupAtmosWinds}
This class is the main class for the model, providing the
external/user-interface into the model.

The model comprises:
\begin{itemize}
\item a set of data tables provided by \textit{DRWPTableLookup}
  instances,
\item the logic for populating, reading, and interpreting a selected table,
  and
\item output variables populated by the results of interpolating the data
  within the selected table.
\end{itemize}
Each table provides an atmospheric profile (as a
function of altitude) for a specific region of the atmosphere. The model
supports mid-run transitions between profiles.

The class data members can be divided into four categories:
\begin{itemize}
\item Model Configuration
\item Table Identification
\item Model Outputs
\item Internal Management
\end{itemize}
We will look at these one at a time.

\subsubsection{LookupAtmosWinds Model Configuration}
The following variables are used for model configuration:

\begin{itemize}
\item \textbf{alt\_bias} is a public variable used to facilitate situations
  in which the altitude provided by the simulation to the
  \textit{LookupAtmosWinds::update( double altitude)} method
  uses a different datum than that used to populate the data files.
\item \textbf{altitude} is the variable used as the independent variable in
  the table interpolation mechanism. This variable is public purely so that it
  is visible to the table-interpolation. Setting this value manually is
  futile. It is computed on each pass of
  \textit{LookupAtmosWinds::update( double altitude)} by combining the altitude
  passed in and
  \textit{alt\_bias} to bias the incoming altitude to account for any offset
  between the datum used for zero-altitude by the simulation, and
  the datum used for zero-altitude by the data tables.
\item \textbf{block\_warnings} is a flag used to block certain messages
  issued by the model, including:
  \begin{itemize}
  \item A warning about the altitude being out of the domain of the current
    data table. This warning would normally be issued once (per table) only
    on the very first occurrence of a domain violation, but can be blocked
    entirely.\\
    \textbf{Note:} Consideration was given to a design in which a first
    domain violation would result in the warning being issued and the flag
    being set to prevent further warnings, then that blocking flag being
    reset if the altitude subsequently returns to within the tabulated domain.
    This would support a warning every time the domain bounds are crossed,
    which would provide, for example, a warning at top and bottom if the
    vehicle transits through the entire domain.
    However, this can also result in
    multiple messages being sent when multiple vehicles are concurrently
    accessing the same atmospheric profile, with one in-domain, and one
    out-of-domain -- one would trigger the warning and set the blocking
    flag, and the next would reset the blocking flag, so the first would
    reissue the warning, etc.
    This consideration was subsequently rejected.
  \item Informational message when loading a new profile mid-run.
  \item Informational message when selecting the initial table by index
    rather than by filename/profile-number.
  \end{itemize}
\item \textbf{gamma} is the ratio of specific heats $C_p / C_v$. For a
  diatomic gas, this has a value of $1.4$; it would be unusual to change
  this value. It is used in the computation of
  the speed of sound from the tabulated atmospheric data.
\item High-altitude speed of sound parameters. At high-altitude, the
  conventional approach for computing speed of sound becomes unreliable. The
  model sets the speed of sound to a constant value above some altitude,
  with a phase-in approach to ensure continuity in the transition below that
  altitude. This
  requires three parameters:
  \begin{itemize}
  \item\textbf{SOS\_hi\_alt\_const} is the high-altitude constant value for
    speed-of-sound, specified in $m s^{-1}$.
  \item\textbf{SOS\_fair\_hi\_alt} is the altitude, specified in $m$,
    above which the
    speed-of-sound is set to the constant value. This also defines the upper
    bound of the region over which the phase-in is applied.
  \item\textbf{SOS\_fair\_lo\_alt} is the lower bound of the
    region over which the phase-in is applied.
  \end{itemize}
\end{itemize}

\subsubsection{LookupAtmosWinds Table Identification}
There are three variables that provide information about the data table.
These variables serve dual purpose depending on context, they may represent
either the table currently being used, or the table about to be used:

\begin{itemize}
\item \textbf{drwpFileName} is the name of the file containing the binary
data used in populating the table.
\item \textbf{wind\_number} is the profile number within the binary file for
the profile used in populating the table. This name is a legacy misnomer,
held over from earlier implementations of this model that were restricted
to providing wind data only; a more appropriate name would be
\textit{profile\_number}.
\item \textbf{include\_vertical\_component} is a flag indicating whether the
table has data for the vertical wind velocity component.
\end{itemize}

When the model selects and loads a new table, these three variables are
assigned the values from the equivalent variables found in the selected
instance of \textit{DRWPTableLookup}.

The model may also use these public variables to identify which table to
select; in this context, they are user-settable values that can be used to
configure the model.  When the \textit{change\_profile()} method is called
(which may be from the \textit{activate()} method as the model is first
activated, or by a deliberate follow-up call once the model is active),
the model will try to find a profile defined by the
specified \textit{drwpFileName} and \textit{wind\_number}. If no profile is
found already loaded, the model will attempt to load one from the specified
file. In this case, the setting of \textit{include\_vertical\_component}
will be used to determine how to interpret the data in the binary file.


\subsubsection{LookupAtmosWinds Model Outputs}

\begin{longtable}{|p{1.5cm}|p{3cm}|p{5.5cm}|p{1cm}|p{2cm}|}
\caption{Atmospheric variables of interest}\\
\hline
\textbf{Type} & \textbf{Name} & \textbf{Purpose} & \textbf{Units} & \textbf{Protection}\\
\hline
\endfirsthead
\multicolumn{5}{c}%
{\tablename\ \thetable\ -- \textit{Continued from previous page}} \\
\hline
\textbf{Type} & \textbf{Name} & \textbf{Purpose} & \textbf{Units} & \textbf{Protection}\\
\hline
\endhead
\hline \multicolumn{5}{r}{\textit{Continued on next page}} \\
\endfoot
\hline
\endlastfoot
double & u & Easterly directed wind velocity & $m\ s^{-1}$ & public \\
\hline
double & v & Topocentric-Northerly directed wind velocity & $m\ s^{-1}$ &
public \\
\hline
double & w & Topocentric-Upwardly directed wind velocity & $m\ s^{-1}$ &
public \\
\hline
double & wind\_angle\_\hspace{0.5cm}blowing\_to & Azimuth angle in the
topocentric-horizontal plane & $rad$ & public \\
\hline
double & wind\_angle\_\hspace{0.5cm}blowing\_from & Azimuth angle in the
topocentric-horizontal plane & $rad$ & public \\
\hline
double & wind\_velocity\_tc & 3-vector for wind velocity in topocentric NED
frame & $m s^{-1}$  $m\ s^{-1}$ & public \\
\hline
double & wind\_vmag & magnitude of wind\_velocity\_tc & $m\ s^{-1}$ & public \\
\hline
double & average\_wind & 3-vector representing an average of
wind\_velocity\_tc over a specified altitude range & $m\ s^{-1}$ & public \\
float & rho & atmospheric density & $kg\ m^{-3}$ & public \\
\hline
float & T & temperature & $K$  & public \\
\hline
float & P & pressure & $Pa$ & public \\
\hline
float & SOS & speed of sound & $m\ s^{-1}$ & public \\
\hline
\end{longtable}

\subsubsection{LookupAtmosWinds Internal Management}
The following data members are used internally only:
\begin{itemize}
\item \textbf{TableLookup\_array} is an array of (10) instances of
\textit{DRWPTableLookup}. This array is sized by:
\item \textbf{MaxDatasets} specifies the number of elements allocated in the
  \textit{TableLookup\_array} array. This is a compile-time setting. It is
  currently set to $10$.
\item \textbf{number\_of\_datasets} records how many of the
  \textit{DRWPLookupTable} instances of \textit{TableLookup\_array} have
  been populated with data.
\item \textbf{current\_index} identifies which of the
  \textit{DRWPLookupTable} instances of \textit{TableLookup\_array} is
  currently being used. This is the index to the array.
\item \textbf{warning\_issued} is a flag indicating that the out-of-domain
  warning has been issued for the current table. This warning is issued no
  more than once per table, being blocked after the first time.
\item \textbf{size\_of\_source\_float} is a compile-time setting, specifying
  how many bytes the floating-point values in the binary data files occupy.
  This is currently set to $4$; the expectation is that the binary data were
  constructed from data of type \textit{float} (and 32-bit integers).
\item \textbf{size\_of\_source\_int} is a compile-time setting, specifying
  how many bytes the integer values in the binary data files occupy.
  This is currently set to $4$; the expectation is that the binary data were
  constructed from 32-bit integers (and data of type \textit{float}).
\end{itemize}


\subsubsection{LookupAtmosWinds::LookupAtmosWinds()}
\textit{Constructor}\\
Creates an instance of the DRWP LookupAtmosWinds class and sets default values.

\subsubsection{LookupAtmosWinds::initialize()}
\textit{Called from: \textbf{external}}\\
Verifies the data content in all profiles that have been
pre-loaded via \textit{load\_DRWP\_file(...)} by calling
\textit{DRWPTableLoookup::verify()} for each pre-loaded member of the
\textit{TableLookup\_array} array.


\subsubsection{LookupAtmosWinds::activate()}
\textit{Called from: \textbf{internal}}\\
Note -- This method is called in response to successful model initialization
and an
active subscription to the model; these can occur in either order. This
pattern adheres to the design of the underlying CML Subscriptions model that
activates a model once it has been initialized and identified as needed with
a \textit{subscribe()} call. The methods
\textit{initialize()} and \textit{subscribe()} are both external calls.

If \textit{drwpFileName} has been specified, activation results in a call to
\textit{change\_profile} to find (or load) a profile that matches the
specification of \textit{drwpFileName} and \textit{wind\_number} (and
\textit{include\_vertical\_component} if loading a new data profile).

If \textit{drwpFileName} has \textbf{not} been specified, activation results
in selecting the profile loaded to the selected index of the
\textit{TableLookup\_array} array. The index can be assigned via the
\textit{change\_datafile\_index} method.

If the selection of data pprofile fails for any reason (e.g. there are no
pre-loaded profiles and no specification of \textit{drwpFileName}, or the
specified profile cannot be found) the activation fails and the model cannot
be used.


\subsubsection{LookupAtmosWinds::update(double altitude)}
\textit{Called from:  \textbf{external}}
\begin{enumerate}
\item Adds an altitude bias (if present) to the input altitude.
\item Calls the \textit{DRWPTableLookup} class's inherited
  \textit{SimpleTableLookup::update()} method for the current data table, which
  performs all of the interpolation.
\item Calls \textit{calculate\_speed\_of\_sound()}
\item Calls \textit{calculate\_wind\_mag\_dir()}
\item Queries the status of the interpolation algorithm to identify whether
  the input altitude is within the tabulated altitude domain. If not, it
  posts a single-issuance warning about being outside the domain.
\end{enumerate}

\textbf{Note:} In many CML models, the \textit{update()} method is
structured to take no arguments. In this case, the model is intended to
simultaneously handle multiple vehicles if necessary, so each vehicle would
need to pass its own altitude in as an argument. The model cannot track
vehicle altitudes because it has no knowledge of which vehicles are using
it. This orgnaization supports the general design paradigm that
\textit{planets} have atmospheres, and \textit{vehicles} use them.


\subsubsection{LookupAtmosWinds::change\_datafile\_index(size\_t)}
\textit{Called from:  \textbf{external, activate(), or change\_profile()}}
This method is used to select which data table to use in subsequent calls
to \textit{update()}. Selection can be applied at any time; if the model is
currently active, this will switch which interpolation-table is currently
being used.

The argument provides the index in the \textit{TableLookup\_array} array that
contains the desired data profile.

If the model is already active, the formerly-current table is deactivated.

The identifying parameters \textit{drwpFileName},  \textit{wind\_number},
and \textit{include\_vertical\_component}, are copied in from their
counterparts in the newly selected table. The \textit{current\_index} is
assigned to the new index, and the new table activated.

\subsubsection{LookupAtmosWinds::change\_profile()}
\textit{Called from:  \textbf{external or activate()}}
Like \textit{change\_datafile\_index(size\_t)}, this method is used to
select which data table to use in subsequent calls
to \textit{update()}. Selection can be applied at any time; if the model is
currently active, this will switch which interpolation-table is currently
being used.

While \textit{change\_datafile\_index(size\_t)} specifies the new table by
index in the \textit{TableLookup\_array} array, this method uses the class
members \textit{drwpFileName}, and \textit{wind\_number} to identify a match
in the \textit{TableLookup\_array} array. If no match is found, the call
forwards to  \textit{load\_DRWP\_file(...)} to find and load a suitable
profile using \textit{drwpFileName}, \textit{wind\_number}, and
\textit{include\_vertical\_component}. If both attempts fail, the method
fails, an error message is posted, and the profile remains unchanged.

\subsubsection{LookupAtmosWinds::load\_DRWP\_file(...)}
\textit{Called from: \textbf{external or change\_profile()}}

The primary purpose of this method is to pre-load data profiles prior to
model initialization. This method can also be accessed internally in
response to initialization without preloading, or switching to a new profile
without preloading. This method performs numerous checks on the integrity
and structure  of
the specified file to be as sure as possible of its correct interpretation.

The method takes three arguments to identify the profile to be loaded:
\begin{enumerate}
\item std::string drwpFileName\_,
\item bool contains\_vertical\_component\_,
\item unsigned int wind\_number\_)
\end{enumerate}

The binary files are assumed to be constructed of combinations of integer
and floating-point values with a specified format and specific data-types.
The model provides for relatively easy adjustment to these default formats,
but those adjustments are made at compile-time, not at run-time.

As discussed previously, the model identifies the expected memory size
occupied by the integer and floating-point values in the binary file with
the class members \textit{size\_of\_source\_integer} and
\textit{size\_of\_source\_float}. These both defalt to $4$ (bytes).

Within this method, two data types are defined, \textit{data\_int}, and
\textit{data\_float}. The data from the binary file will be read into
individual instances and arrays of instances of these variables. These
are currently implemented as \textit{int32\_t} and \textit{float}
respectively (both of which are 32-bit / 4-byte). The first
check the method makes is that these two data types have sizes that match
the expected size in the binary data.

The method then checks for the availability of at least one slot in the
\textit{TableLookup\_array} array. This array has a fixed (compile-time)
size, so if all slots are already populated, we cannot add a new data
profile.

Then the binary file is opened and systematically parsed and read, checking
for expected values wherever possible to verify the integrity and structure
of the file. The expected structure of the file is as follows:
\begin{enumerate}\label{sec:binary_read_struc}
\item An integer specifying the number of profiles in the file
  (\textit{nProfiles}).
\item An integer specifying the number of altitude breakpoints /
  calibration-points used for each profile (\textit{nAlts}). Note -- all
  profiles within a file use the same set of altitude breakpoints /
  calibration-points.
\item An array of \textit{nAlts} floating-point values specifying the values
  of the altitude breakpoints / calibration-points.
\item An array of \textit{nProfiles} integer values specifying the profile
  numbers of each profile in the order in which they are included in the
  binary file.
\item An array of \textit{nProfiles} profiles. Each profile is defined by
the following data:
  \begin{enumerate}
  \item An integer defining the profile number of this profile. This must
    match with the appropriate value in the earlier array of profile
    numbers.
  \item An array of \textit{nAlts} floating-point values specifying the
    values for the West-to-East wind component.
  \item An array of \textit{nAlts} floating-point values specifying the
    values for the South-to-North wind component.
  \item \textbf{(optional)}An array of \textit{nAlts} floating-point
    values specifying the
    values for the Upward-positive vertical wind component.
  \item An array of \textit{nAlts} floating-point values specifying the
    values for the temperature.
  \item An array of \textit{nAlts} floating-point values specifying the
    values for the density.
  \item An array of \textit{nAlts} floating-point values specifying the
    values for the pressure.
  \end{enumerate}
\end{enumerate}

For each of the following reads from the file, if the read itself fails
(likely due
to reading past the end of the file), or the read result is inconsistent with
expectation, the method is terminated. This usually results in simulation
termination.

The first and second integers are read and compared against reasonable
expectation. If either
value is 0 (indicates empty file) or very large (typically indicative of
an unmatched memory size between the way the file was written and the way it
ws read), the overall method is terminated.

The altitude array is read next, and if successful, its values assigned to a
STL-vector of type \textit{double}. This vector is passed to the CML
table-interpolation model for loading into the interpolator's independent
variable data set. This action includes checks on things like monotonicity
of the altitude data. If the loading of data to the interpolator fails, the
overall method is terminated.

The array of profile numbers is read next. There is no independent
evaluation that can be applied to these values but once read, the array is
searched for the presence of the desired profile number, and the index of
that match recorded. If no match is found, the overall method is terminated.

At this point, we do not need to read the entirety of the rest of the file
containing all of
the profiles, we only need the one desired profile from this file,
and we know where that should be
located. To make sure that the rest of the file is structured as expected
before jumping to the desired profile start, we make two additional sanity
checks:

\begin{itemize}
\item
The next field is read as an integer, this should be the profile number of
the first profile in the list, and should match the first entry of the array
of profile numbers already read in. If it does not, the method is
terminated. We now skip over enough memory to house the floating-point
values of the profile (which is a function of the number of altitude
breakpoints, and whether the profiles include the vertical velocity
component) and read the next field as an integer, which should match with
the second profile number. If it does not, the method is terminated. This
process is repeated for all \textit{nProfiles} profiles.
\item
After reaching the end of the last profile, the position is compared
against the size of the file to make sure the file is appropriately sized
to house the specified number of profiles. If this does not match, the
method is terminated.
\end{itemize}

Now, with confidence that the file structure matches with expectation,
we return to the
position at the start of the profile of interest. We ignore the first field
for the profile number (that value has already been checked), and read in
sufficient
data to cover all the variables for each altitude breakpoint. As with the
altitude data, these data are assigned to a
STL-vector of type \textit{double}, in preparation for the CML
table-interpolation model's data loading.

Next, we populate a STL-vector of pointers with the addresses of the class
members \textit{u, v, w, T, rho, P} (\textit{w} being optional). The data vector
and the variables vector are both passed to the CML table-interpolation
model's method for loading onto the interpolator's
\textit{dependent} variables data set. If this data-loading process fails,
the method is terminated.

At this point, the CML table-interpolation model has its data tables
fully populated and can be used to interpolate between the specified data
points. The binary file can be closed.


\subsubsection{LookupAtmosWinds::compute\_average\_wind(...)}
\textit{Called from: \textbf{external}}

This method computes the average of the wind velocity vectors, outputting
the computed average into the 3-element vector \textit{average\_wind} with
the three components representing average expressed in the topocentric
NED frame.

This method has four different calling interfaces with 0, 1, 2, or 3
arguments. We will look at the 3-argument method, the others are derived
from it.
\begin{enumerate}
\item \textbf{table index} The averages will be taken from the data table at
  the specified index in the \textit{TableLookup\_array} array. If this is
  not specified (0 or 2 argument signatures):
  \begin{itemize}
  \item if the model is active, it will use the table currently in use.
  \item if the model is not active, it will the first table in the
    \textit{TableLookup\_array} array.
  \end{itemize}
\item \textbf{minimum altitude}
\item \textbf{maximum altitude} The minimum and maximum altitudes are
  presented as a pair of arguments. The model supports specification of an
  altitude domain over which to compute the average winds.
  If these are omitted (0 or 1 argument
  signatures), the full domain is used for the averaging.
\end{enumerate}

Behind the scenes, the four interface methods all feed into a single,
4-argument method. This method is protected so cannot be externally called.
The hidden method proceeds as follows:

\begin{enumerate}
\item The selected index is compared against \textit{number\_of\_datasets}
  to confirm that this table has been populated. If it has not, the method
  returns with no action.
\item The set of altitude values in the lookup-table's independent variable
  array is stepped
  through sequentially. Each altitude point is used if either the method is
  using the full domain, or the point value falls between the minimum and
  maximum altitudes specifying the domain over which to compute the average.
  For all altitude points that meet this criterion, the model variable
  \textit{altitude} is
  assigned to this calibration point, and the interpolation table updated,
  assigning values to \textit{u, v, w}. The
  \textit{calculate\_wind\_mag\_dir()} converts these values to the
  topocentric-NED 3-element vector, which is then be incremented onto a
  cumulative total.
\item The accumulated total of the 3-element vectors is then divided by the
  number of points used to give a mean average.
\item Finally, to ensure that the model outputs are consistent with the
  model's actual intent (to compute winds and atmosphere at a vehicle's
  altitude), the model's \textit{update()} method is executed again to
  recompute the state based on the actual vehicle altitude, returning
  the model's output variables to their original state.
\end{enumerate}

\textbf{Note:}
There is significant ambiguity in what an ``average velocity'' means. This
implementation is an average over the calibrated points, but this is not the
same as a spatial average because there is no verification that the data
points are equally distributed. A spatially averaged wind velocity would
have to weight each data point by the spatial separation between it and its
neighbors.\\
Neither does this represent a temporally-averaged velocity, such as
would be encountered by a vehicle on
some trajectory; this would also require the local vehicle velocity to
support weighting each calibrated value by the duration the vehicle spends
in its vicinity.\\
Furthermore, if the desired representation is that of a wind representing
the average dynamic effect on a vehicle, this average falls woefully short.
Such an average would require
squaring the relative speed along the direction of interest and multiplying
by both the duration for which that speed is relevant and the density of the
atmosphere in the vicinity, even if we assume simple aerodynamics such as a
coefficient-of-drag or ballistic-coefficient type representation.\\
Simply averaging the wind velocities provides
information, but it is highly subject to misinterpretation and
misapplication.



\subsubsection{LookupAtmosWinds::(calculate\_speed\_of\_sound)}
\textit{Called from:  \textbf{update()}}

See \ref{sec:sos_algo}

\subsubsection{calculate\_wind\_mag\_dir()}
\textit{Called from:  \textbf{update(), compute\_average\_wind(...)}}

See \ref{sec:wind_dir_mag}.

\subsubsection{LookupAtmosWinds::stream\_error(...)}
\textit{Called from:  \textbf{load\_DRWP\_file(...)}}

Provides shared error-handling messages for multiple potential errors coming
from the \textit{load\_DRWP\_file(...)} method.

\subsubsection{LookupAtmosWinds::set\_include\_vertical\_component(bool val)}
\textit{Called from:  \textbf{deprecated}}

This method was deprecated in the 2023 refactor because its old capability is
no longer needed. However, it was accessible to the public interface, so it has
been replaced with a method that simply issues an error message.

\subsubsection{LookupAtmosWinds::test\_for\_reinitialize()}
\textit{Called from:  \textbf{deprecated}}

This method was deprecated in the 2023 refactor because its old capability is
no longer needed. However, it was accessible to the public interface, so it has
been replaced with a method that simply issues an error message.



\section{Mathematical Formulation}
\subsection{Data requirements}

The DRWP model requires that the independent variable (the altitude) be
everywhere monotonic.

\subsection{Algorithms}
\subsubsection{One-dimensional interpolation}
The interpolation is handled by the CML table-interpolation method. See
model documents for details
(\textit{cml/models/utilities/table\_interp\_cpp/docs/table\_interp.pdf)}.

General principle of linear interpolation:
Let $x$ be the altitude of interest, and $x_{n}$, $x_{n+1}$ be the
bracketing values of the this independent variable, such that
$x_{n}$ $\leq$ $x$ $\leq$ $x_{n+1}$ or
$x_{n+1}$ $\leq$ $x$ $\leq$ $x_{n}$.  Then define a value to represent the
fractional part of that interval:
\begin{equation}
  \alpha = \frac{x -  x_{n}} {x_{n+1} - x_{n}}
\end{equation}
The desired interpolated value, $f(x)$, then becomes
\begin{equation}\label{eq:interp}
  f(x) = f(x_{n}) + \alpha \cdot (f(x_{n+1}) - f(x_{n}))
\end{equation}

\subsubsection{Speed of Sound}\label{sec:sos_algo}
If we let $\gamma = \frac{C_{p}}{C_{v}}$, where $C_{p}$, $C_{v}$ are the
heat capacities at constant pressure and constant volume, respectively.
Then the speed of sound $c$ is obtained from
\begin{equation}\label{eq:SOS}
  c = \sqrt{\gamma\frac{P}{\rho}}
\end{equation}
where:\\
\hspace{1cm}$P$ = pressure at altitude,\\
\hspace{1cm}$\rho$ = density at altitude

Standard atmospheric conditions give $\gamma$ a value of 1.4 (diatomic gas
$N_2$ AND $O_2$).

At high altitudes, an adjustment must be made to equation \ref{eq:SOS}.
$c$ is considered to remain constant at high altitudes, so a transition
(fairing) must occur from equation \ref{eq:SOS} to this altitude.
This transition uses a rolling interpolation of the form:
\begin{equation}\label{eq:SOS_fair}
  c = \sqrt{\gamma\frac{P}{\rho}} + \frac {h - h_{lo}}{h_{hi} - h_{lo}} \cdot
  \left( c_{hi} -  \sqrt{\gamma\frac{P}{\rho}} \right)
\end{equation}
where:\\
%\begin{adjustwidth}{2.5em}{0pt}
~~$h$ = altitude of interest,\\
~~$h_{lo}$ is the altitude at the floor of the transition region\\
~~$h_{hi}$ is the altitude at the ceiling of the transition region\\
~~$c_{hi}$ is the speed of sound at the ceiling of the transition region\\
%\end{adjustwidth}

Note the apparent similarity but fundamental difference between equation
\ref{eq:SOS_fair} and \ref{eq:interp}. This is not a linear interpolation
between the speeds of sound at the floor and ceiling of the transition
region. It is a linear interpolation between the speeds of sound at the
\textbf{current altitude} and ceiling of the transition region,
with the relative weightings determined by the relative position between
the \textbf{floor} and ceiling of the transition region. This produces a
profile that generally follows the low-altitude equation, gradually diverging
from it until it smoothly reaches $c_{hi}$ at $h_{hi}$.

Default values for this transition region are:
\begin{itemize}
\item $h_{lo} = 85,344 m$
\item $h_{hi} = 91,440 m$
\item $c_{hi} = 274.61 m\ s^{-1}$
\end{itemize}

For altitudes above $h_{hi}$,
\begin{equation}
  c = c_{hi}
\end{equation}

\clearpage
\subsubsection{Wind magnitude and direction}\label{sec:wind_dir_mag}
DRWP calculates the wind direction and magnitude by
\begin{equation}
\theta = tan^{-1}\left(\frac{u}{v}\right)
\end{equation} and
\begin{equation}
| \vec v| = \sqrt{u^2 + v^2 + w^2}
\end{equation}
where:
$u, v, w$ are the East, North, and Up components of the
wind velocity, respectively.

These values are available as the public members
\textit{wind\_angle\_blowing\_to} and \textit{wind\_vmag}.


\chapter{User's Guide}

\section{Instantiation and Construction}
The DRWP model is typically instantiated as part of the Earth planet
object. A single instance of this model can service multiple vehicles
simultaneously.

\begin{code}
LookupAtmosWinds   DRWP;
\end{code}

The model constructor takes no arguments.


\section{Implementation via Atmosphere-Executive Interface}
It is recommended to use the CML \textit{Atmosphere Executive}
model (see \textit{cml/models/environment/atmos/atmos\_exec} for managing
the implementation of this model within a simulation framework. The
atmosphere-executive model handles switching between atmospheres; it is
equipped with the code to manage the associated initialization, activation,
execution, and deactivation processes of the DRWP model for any vehicle
electing to use DRWP for representation of the atmosphere in its locality.
See the Atmosphere-executive model documentation for
details on instantiating an Atmosphere-executive, and configuring it to use
the DRWP model.

\section{Standalone Implementation}
It is also possible to use the DRWP model independently of the
Atmosphere-Executive. In this case, there are four methods that the user
must be familiar with.
\begin{itemize}
\item \textbf{initialize()} should be executed during simulation
  initialization. The model cannot be activated without having executed this
  method.
\item \textbf{subscribe()} is executed typically by a vehicle object when it
  starts to need atmospheric data in its locality. If this is the first (or
  only) subscription to the model, this call will activate the model.
\item \textbf{unsubscribe()} is executed when the vehicle object no longer
  needs an updated atmospheric conditions. This removes a subscription. If
  the last subscription is removed, the model deactivates and subsequent
  calls to \textit{update( altitude)} will return without action.
\item \textbf{update( double altitude)} is typically executed by the vehicle
  object, passing in the value of the vehicle's current altitude. This call
  may be scheduled for regular execution even while the model is not needed;
  if the model is not subscribed, this will have no effect. However, it is
  more efficient to call this method only as needed.
\end{itemize}

\section{Model Data}
\subsection{Default values}
Default data is all provided within the model constructor; there is no
separate method to populate the model with default values.
Most values are constructed to 0, 0.0, or false (depending on the data type).
Default values used in the speed of sound calculations are presented in
section \ref{sec:sos_algo}.

\section{Setting Input Data}\label{sec:input_data}
The user must supply a file with the profile data.
This file must be in binary format with a structure presented in section
\ref{sec:binary_read_struc}. That structure is summarized here for
convenience.

\begin{itemize}
\item A 32-bit integer representing the number of profiles
\item A 32-bit integer representing the number of altitudes per profile
\item Array of \textit{float} values representing the altitudes
  (the independent variable)
\item Array of 32-bit integers representing the profile-numbers for
  profiles included in this file.
\item Array of compound data representing the data for each profile. 
  For each profile:
  \begin{itemize}
  \item A 32-bit integer representing the specific profile number
  \item Array of \textit{float} values representing the U wind velocity
    component values (1 value for each altitude in the altitude-array).
  \item Array of \textit{float} values representing the V wind velocity
    component values (1 value for each altitude in the altitude-array).
  \item (optional) Array of \textit{float} values representing the W
    wind velocity component values (1 value for each altitude in the
    altitude-array).
  \item Array of \textit{float} values representing the temperature
    (1 value for each altitude in the altitude-array).
  \item Array of \textit{float} values representing the density
    (1 value for each altitude in the altitude-array).
  \item Array of \textit{float} values representing the pressure
    (1 value for each altitude in the altitude-array).
  \end{itemize}
\end{itemize}

The user must also supply the filename of the file containing the binary
file, and the profile number to find the desired profile within said file.
The profile number and name of the binary file are typically
supplied in the simulation input file.

\textbf{Notes:}
\begin{itemize}
\item If the configuration specifies inclusion of the vertical
  component of wind velocity, it must be present (can be zero) in all
  profiles in the specified data file.
  Similarly, if the vertical component of wind velocity is present (can be
  zero) in the data files, it must be specified for inclusion in the
  interpolation tables.
\item All floating-point values (altitude and atmospheric variables) must be
  of the same type, as must all
  integer values (sizes and indices). By default, the model expects these to
  be 32-bit (\textbf{float} and \textbf{int32\_t}).
\end{itemize}


\section{Extracting Output Data}
The model outputs include the three main wind components,
wind magnitude and direction, density, temperature, pressure,
and the speed of sound.  These are available as public members of
the instantiated class as data type \textit{double}.
The most commonly used are:

\begin{itemize}
\item u
\item v
\item w
\item wind\_velocity\_tc (3-element array)
\item wind\_angle\_blowing\_to
\item wind\_angle\_blowing\_from
\item wind\_vmag
\item rho
\item T
\item P
\item SOS
\item average\_wind (3-element array)
\end{itemize}




%***********************************************************************
% Verification chapter
%***********************************************************************
\chapter{Verification}
The model was inspected to ensure compliance with coding standards.
Verification of results was accomplished through unit simulation testing.
Test altitudes are managed by a sweep algorithm in the Unit-test-framework
model, sweeping from $-3\ km$ to $+187\ km$ in $2\ km$ intervals.
The binary data files are found in \textit{SIM\_lookup/Binaries}. This
directory also includes csv files corresponding to selected profiles from the
binary data for purposes of comparing the model outputs against a
human-readable data set.

The swept altitude domain exceeds the domain of the binary data files on both
ends; this will be identified and noted by the model upon execution.

\section{Code Coverage}
\begin{figure}[h!]
   \centering
   \includegraphics[width=\linewidth]{Images/code_coverage.png}
\end{figure}

\subsection{Exceptions}
\subsubsection{lookup\_winds.hh}
The header file has no methods to cover; the two lines identified and not
covered are the empty class destructors.

\subsubsection{lookup\_winds.cc}
There are six lines not covered in this file. Each of these lines would be 
executed upon detection of one of five errors. Each error:
\begin{itemize}
\item appears to be  unreachable through the automated unit-testing
process,
\item has been commented as such in the code, and
\item has been tested manually in gdb to execute the lines not covered in
automated testing.
\end{itemize}


\section{Input data}
Some verification tests (such as \textit{RUN\_01\_preload\_5\_nospec}) read
the binary file\\
\textit{Binaries/DRWP\_no\_w\_comp.bin}.
This binary file does not provide the vertical velocity component data so the
interpolation populates five dependent variables.

Some verification tests (such as \textit{RUN\_02\_preload\_6\_nospec}) read
the binary file\\
\textit{Binaries/DRWP\_w\_comp.bin}.
This binary file does provide the vertical velocity component data so the
interpolation populates six dependent variables.

In \textit{verif/SIM\_lookup/Binaries/corrupt\_binaries}, a set of partial
binary files exist to test the error detection in the binary file reader.

Each of the two valid binary files contains multiple profiles from which one
profile is used repeatedly in the unit-tests discussed below.
Human-readable, comma-separated versions of these two frequently-used
profiles are also made available in the \textit{Binaries} folder\\
(\textit{Binaries/DRWP\_no\_w\_comp\_\_prof\_1771.csv} and
\textit{Binaries/DRWP\_w\_comp\_\_prof\_2.csv}).

\section{Results}
For each verification run, the configuration and purpose of the run is
listed, followed by the results of the run.

%***********************************************************************
% RUN_01
%***********************************************************************
\subsection{RUN\_01\_preload\_5\_nospec}\label{sec:ver:RUN01}
\textbf{Configuration:}\\ preload\\
\textit{Binaries/DRWP\_no\_w\_comp.bin}, profile $1771$

\textbf{Purpose:}\\ To test the interpolation of a 5-element data file.

This scenario uses the \textit{load\_DRWP\_file(...)} method to preload
the specified profile. At execution, the model defaults to using the
preloaded profile from index 0.

The model correctly identified out-of-range input altitudes on the low side
of the altitude domain.
Only the first occurrence outside the
domain generated a warning, subsequent warnings were inhibited.

\textbf{Note:} Because this warning is only issued once, it is not
repeated when the altitude goes out of domain on the high side.
\begin{figure}[h!]
  \includegraphics[width=\linewidth]{Images/OutOfRange.png}
\end{figure}

The wind direction and magnitude were calculated for each altitude.
See \ref{sec:wind_dir_mag}.

The profile used for this test has calibrated data points at convenient
round-number values of altitude. The values of altitude used as inputs are
found in the data tables so the values for atmospheric variables returned by
the model should be identical to those found in the data table, an example
is shown in Figure \ref{fig:01_9}.

\begin{figure}[h!]
  \centering
  \includegraphics[width=0.45\linewidth]{Images/RUN_01_figure_9.png}
  \caption{For this profile, all lookup-altitudes overlap with
    calibration-points.}
  \label{fig:01_9}
\end{figure}

The raw data and interpolated data are compared for all variables in the
Figures \ref{fig:RUN01_1} - \ref{fig:RUN01_3}.

\begin{figure}[ht!]
  \centering
  \begin{subfigure}[b]{0.45\linewidth}
    \includegraphics[width=\linewidth]{Images/RUN_01_figure_1.png}
  \end{subfigure}
  \begin{subfigure}[b]{0.45\linewidth}
    \includegraphics[width=\linewidth]{Images/RUN_01_figure_2.png}
  \end{subfigure}
  \caption{Comparison of Wind Velocity components, including raw data,
  interpolated data (\textit{u,v}), and NED output data
  (\textit{wind\_velocity\_tc}).}
  \label{fig:RUN01_1}
\end{figure}

\begin{figure}[ht!]
  \centering
  \begin{subfigure}[b]{0.45\linewidth}
     \includegraphics[width=\linewidth]{Images/RUN_01_figure_4.png}
  \end{subfigure}
  \begin{subfigure}[b]{0.45\linewidth}
     \includegraphics[width=\linewidth]{Images/RUN_01_figure_5.png}
  \end{subfigure}
  \begin{subfigure}[b]{0.45\linewidth}
    \includegraphics[width=\linewidth]{Images/RUN_01_figure_6.png}
  \end{subfigure}
  \caption{Comparison of density, pressure, and temperature profiles from
  raw data and interpolated output.}
  \label{fig:RUN01_3}
\end{figure}

\newpage

The profiles of the derived quantities, speed of sound and wind speed,
are shown in Figure \ref{fig:RUN01_4}. These values are calculated from the
interpolated values at each test altitude; these are not provided in raw
data. See sections \ref{sec:sos_algo} and \ref{sec:wind_dir_mag}.

\begin{figure}[h!]
  \centering
  \begin{subfigure}[b]{0.45\linewidth}
    \includegraphics[width=\linewidth]{Images/RUN_01_figure_7.png}
  \end{subfigure}
  \begin{subfigure}[b]{0.45\linewidth}
    \includegraphics[width=\linewidth]{Images/RUN_01_figure_8.png}
  \end{subfigure}
  \caption{Speed of sound and wind speed computed profiles}
  \label{fig:RUN01_4}
\end{figure}

\clearpage
\newpage

%***********************************************************************
% RUN_02
%***********************************************************************
\subsection{RUN\_02\_preload\_6\_nospec}\label{sec:ver:RUN02}
\textbf{Configuration:}\\ preload\\
\textit{Binaries/DRWP\_w\_comp.bin}, profile $2$

\textbf{Purpose:}\\ To test the interpolation of a 6-element data file.

This scenario uses the \textit{load\_DRWP\_file(...)} method to 
preload the specified profile.
At execution, the model defaults to using the preloaded profile from index 0.

The profile used for this test has calibrated data points at obscure
altitudes that do not match with the round-number values of altitude
provided by the unit test. The values of altitude used as inputs are
\textbf{not} found in the data tables so the values for atmospheric
variables returned by the model should require interpolation between the
values found in the data table, as illustrated in figure \ref{fig:02_9}.
\begin{figure}[h!]
  \centering
  \includegraphics[width=0.45\linewidth]{Images/RUN_02_figure_9.png}
  \caption{For this profile, all lookup-altitudes require innterpolation
    between the calibration-points.}
  \label{fig:02_9}
\end{figure}

As with \textit{RUN\_01}, the model correctly identified out-of-domain
input altitudes on the low sides of the altitude domain, with only the first
occurrence on each side generating a warning and the remaining incidences of
out-of-domain inputs being processed without warnings issued.

The model correctly performed a 1D interpolation on each of the six
atmospheric properties for the given altitude as shown in the Figures
\ref{fig:RUN02_1} - \ref{fig:RUN02_2}.

\begin{figure}[h!]
  \centering
  \begin{subfigure}[b]{0.45\linewidth}
     \includegraphics[width=\linewidth]{Images/RUN_02_figure_1.png}
  \end{subfigure}
  \begin{subfigure}[b]{0.45\linewidth}
     \includegraphics[width=\linewidth]{Images/RUN_02_figure_2.png}
  \end{subfigure}
  \begin{subfigure}[b]{0.45\linewidth}
     \includegraphics[width=\linewidth]{Images/RUN_02_figure_3.png}
  \end{subfigure}
  \caption{Comparison of Wind Velocity components, including raw data,
  interpolated data (\textit{+u,+v,-w}), and NED output data
  (\textit{wind\_velocity\_tc}).}
  \label{fig:RUN02_1}
\end{figure}


\begin{figure}[h!]
  \centering
  \begin{subfigure}[b]{0.45\linewidth}
    \includegraphics[width=\linewidth]{Images/RUN_02_figure_4.png}
  \end{subfigure}
  \begin{subfigure}[b]{0.45\linewidth}
    \includegraphics[width=\linewidth]{Images/RUN_02_figure_5.png}
  \end{subfigure}
  \begin{subfigure}[b]{0.45\linewidth}
    \includegraphics[width=\linewidth]{Images/RUN_02_figure_6.png}
  \end{subfigure}
  \caption{Comparison of density, pressure, and temperature profiles from
  raw data and interpolated output.}
  \label{fig:RUN02_2}
\end{figure}

\newpage

The profiles of the derived quantities, speed of sound and wind speed,
are shown in Figure \ref{fig:RUN02_3}. These values are calculated from the
interpolated values at each test altitude; these are not provided in raw
data. See sections \ref{sec:sos_algo} and \ref{sec:wind_dir_mag}.

\begin{figure}[ht!]
  \centering
  \begin{subfigure}[b]{0.45\linewidth}
    \includegraphics[width=\linewidth]{Images/RUN_02_figure_7.png}
  \end{subfigure}
  \begin{subfigure}[b]{0.45\linewidth}
    \includegraphics[width=\linewidth]{Images/RUN_02_figure_8.png}
  \end{subfigure}
  \caption{Speed of sound and wind speed computed profiles}
  \label{fig:RUN02_3}
\end{figure}

\clearpage
\newpage


%***********************************************************************
% RUN_03-06
%***********************************************************************
\subsection{RUN\_03\_preload\_56\_nospec}
\textbf{Configuration:}\\ preload\\
\textit{Binaries/DRWP\_no\_w\_comp.bin}, profile $1771$ (\textit{RUN\_01})
\textit{Binaries/DRWP\_w\_comp.bin}, profile $2$ (\textit{RUN\_02}).

\textbf{Purpose:}\\ To test the default behavior of defaulting to index 0 when
multiple profiles are pre-loaded.

This scenario uses the \textit{load\_DRWP\_file(...)} method 
to preload the specified
profiles. At execution, the model defaults to using the preloaded profile from
index 0. The results match with those from \textit{RUN\_01} (section
\ref{sec:ver:RUN01}).


\subsection{RUN\_04\_preload\_56\_spec\_6\_file}
\textbf{Configuration:}\\ preload\\
\textit{Binaries/DRWP\_no\_w\_comp.bin}, profile $1771$ (\textit{RUN\_01})
\textit{Binaries/DRWP\_w\_comp.bin}, profile $2$ (\textit{RUN\_02}).

\textbf{Purpose:}\\ To test the ability to select between multiple pre-loaded
profiles using the \textit{drwpFileName} and \textit{wind\_number} settings.

This scenario uses \textit{RUN\_03} as a baseline to preload 2 profiles.
It also sets the \textit{drwpFileName} and \textit{wind\_number}
variables to \textit{Binaries/DRWP\_w\_comp.bin} and value $2$ respectively.
At execution, the model identifies the profile at index 1 as matching the
specification, and runs with that profile.
The results match with those from \textit{RUN\_02} (section
\ref{sec:ver:RUN02}).


\subsection{RUN\_05\_preload\_56\_spec\_6\_ix}
\textbf{Configuration:}\\ preload\\
\textit{Binaries/DRWP\_no\_w\_comp.bin}, profile $1771$ (\textit{RUN\_01})
\textit{Binaries/DRWP\_w\_comp.bin}, profile $2$ (\textit{RUN\_02}).

\textbf{Purpose:}\\ To test the ability to select between multiple pre-loaded
profiles by specifying the desired index.

This scenario uses \textit{RUN\_03} as a baseline to preload 2 profiles.
It also sets desired index to $1$ using \textit{change\_datafile\_index(1)}.
The results match with those from \textit{RUN\_02} (section
\ref{sec:ver:RUN02}).


\subsection{RUN\_06\_spec\_6\_file}
\textbf{Configuration:}\\ preload none\\
run-time load\\
\textit{Binaries/DRWP\_w\_comp.bin}, profile $2$ (\textit{RUN\_02}).

\textbf{Purpose:}\\ To test the ability to identify and load a profile at model
activation when there are no profiles pre-loaded.

This scenario uses no preloads of data profiles.
It sets the \textit{drwpFileName} and \textit{wind\_number}
variables to \textit{Binaries/DRWP\_w\_comp.bin} and $2$ respectively.
At execution, the model identifies that this profile has not been pre-loaded,
and loads it.
The results match with those from \textit{RUN\_02} (section
\ref{sec:ver:RUN02}).


%***********************************************************************
% RUN_07
%***********************************************************************
\subsection{RUN\_07\_mid\_run\_switch\_index} \label{sec:verif:RUN07}
\textbf{Configuration:}\\ preload\\
\textit{Binaries/DRWP\_w\_comp.bin}, profile $2$ (\textit{RUN\_02})\\
\textit{Binaries/DRWP\_w\_comp.bin}, profile $3$\\
\textit{Binaries/DRWP\_no\_w\_comp.bin}, profile $1771$ (\textit{RUN\_01})

\textbf{Purpose:}\\ To test the ability to switch between profiles
mid-run using \textit{change\_datafile\_index(...)}.

The run uses \textit{change\_datafile\_index(...)} to switch between
profiles mid-run. The profile selection cycles from
index 0 to 1 to 2 and back to 0.
Initially, the data should match with \textit{RUN\_02} for ten cycles, then
switch to new data for ten cycles as we use a new profile, then switch to
match with output from \textit{RUN\_01} for ten cycles, and finally switch
back to match with data from \textit{RUN\_02}. In Figure \ref{fig:RUN07_1},
we see the red line from this run overlap the blue line from
\textit{RUN\_02} and the black line from \textit{RUN\_01} as expected.

\begin{figure}[ht!]
  \centering
  \begin{subfigure}[b]{0.45\linewidth}
     \includegraphics[width=\linewidth]{Images/RUN_07_figure_1.png}
  \end{subfigure}
  \begin{subfigure}[b]{0.45\linewidth}
     \includegraphics[width=\linewidth]{Images/RUN_07_figure_2.png}
  \end{subfigure}
  \begin{subfigure}[b]{0.45\linewidth}
     \includegraphics[width=\linewidth]{Images/RUN_07_figure_3.png}
  \end{subfigure}
  \caption{Comparison of Wind Velocity components from \textit{RUN\_01},
    \textit{RUN\_02}, and \textit{RUN\_07}.}
  \label{fig:RUN07_1}
\end{figure}

\clearpage


%***********************************************************************
% RUN_08
%***********************************************************************
\subsection{RUN\_08\_mid\_run\_switch\_spec}
\textbf{Configuration:}\\ preload\\
\textit{Binaries/DRWP\_w\_comp.bin}, profile $2$ (\textit{RUN\_02})\\
\textit{Binaries/DRWP\_w\_comp.bin}, profile $3$\\
\textit{Binaries/DRWP\_no\_w\_comp.bin}, profile $1771$ (\textit{RUN\_01});\\
run-time load\\
\textit{Binaries/DRWP\_w\_comp.bin}, profile $1$

\textbf{Purpose:}\\ To test the ability to switch between profiles --
pre-loaded and unloaded -- mid-run using
\textit{change\_profile()}. The run cycles from
\begin{enumerate}
\item t=0: index 0: by default
\item t=10: index 1: \textit{wind\_number} is set to $3$ and
\textit{drwpFileName} remains unchanged. \textit{change\_profile()}
recognizes these settings as consistent with index 1, and and switches the
selection of data-table.
\item t=20: index 2: \textit{wind\_number} is set to $1771$ and
\textit{drwpFileName} to \textit{Binaries/DRWP\_no\_w\_comp.bin}.
\textit{change\_profile()}
recognizes these settings as consistent with index 2, and switches the
selection of data-table.
\item t=30: index 0: \textit{wind\_number} is set to $2$ and
\textit{drwpFileName} to \textit{Binaries/DRWP\_w\_comp.bin}.
\textit{change\_profile()}
recognizes these settings as consistent with index 0, and switches the
selection of data-table.
\item t=40: index 0: \textit{wind\_number} is set to $1$ and
\textit{drwpFileName} remains unchanged.
\textit{change\_profile()}
recognizes these settings as inconsistent with anything currently loaded,
identifies a matching profile in the specified file, loads it, and
switches the selection of data-table to index 3.
\end{enumerate}

For the first 39 cycles, the data match with those from
\textit{RUN\_07} (section \ref{sec:verif:RUN07}).  The cycles following
where \textit{RUN\_07} stops follows a new profile.

\begin{figure}[ht!]
  \centering
  \begin{subfigure}[b]{0.45\linewidth}
     \includegraphics[width=\linewidth]{Images/RUN_08_figure_1.png}
  \end{subfigure}
  \begin{subfigure}[b]{0.45\linewidth}
     \includegraphics[width=\linewidth]{Images/RUN_08_figure_2.png}
  \end{subfigure}
  \begin{subfigure}[b]{0.45\linewidth}
     \includegraphics[width=\linewidth]{Images/RUN_08_figure_3.png}
  \end{subfigure}
  \caption{Comparison of Wind Velocity components from \textit{RUN\_01},
    \textit{RUN\_02}, and \textit{RUN\_08}.}
  \label{fig:RUN08_1}
\end{figure}

\clearpage
%***********************************************************************
% RUN_10
%***********************************************************************
\subsection{RUN\_10\_alt\_bias}
\textbf{Configuration:}\\ preload\\
\textit{Binaries/DRWP\_no\_w\_comp.bin}, profile $1771$ (\textit{RUN\_01})

\textbf{Purpose:}\\ To test the application of an altitude bias.

This run uses \textit{RUN\_01} as a baseline and applies an bias of +4000
m to the input altitude using the \textit{alt\_bias} variable. This is an
abbreviated run, the intent is to confirm the offset and this can be
accomplished with a handful of data points.

With a 4000 m bias, the unit-test altitude profile starts above the floor
of the model data so the out-of-domain warning message issued in 
\textit{RUN\_01} is not issued for this run.
If the run was allowed to proceed through the full $95$ cycles examined in 
\textit{RUN\_01} (comment out the stop time in the input file),
the out-of-domain warning is issued when the altitude
passes through the domain ceiling.
\includegraphics[width=\linewidth]{Images/RUN_10_warning.png}

With each data point covering a height of 2000 m, the 4000 m bias should
produce an offset in the data logging of two steps (two lines) from the data
produced by \textit{RUN\_01} and should be otherwise identical to it. This
is observed in the output data files
\begin{figure}[h!]
\centering
\includegraphics[width=\linewidth]{Images/data_01_10.png}
\includegraphics[width=\linewidth]{Images/data_10.png}
\caption {Data for tests 2-6 from \textit{RUN\_01},
          and tests 0-4 from \textit{RUN\_10}.
          The altitudes for these two test-sets match and the consequential
          output data also match.}
\end{figure}

The profile resulting from the application of this bias results in a shift
of the existing profile to lower altitudes, as shown in Figure
\ref{fig:RUN10_2}. It may be initially counterintuitive as to why the data
shifts to lower altitude with the application of a positive bias, but this
is the correct behavior. The altitude as provided to the data-lookup process
(the value plotted on the x-axis in Figure \textit{RUN\_07}) is
biased positively, so obtains a value (red line) that would otherwise have been
obtained from a higher altitude (black line).

\begin{figure}[ht!]
  \centering
  \begin{subfigure}[b]{0.45\linewidth}
     \includegraphics[width=\linewidth]{Images/RUN_10_figure_1.png}
  \end{subfigure}
  \begin{subfigure}[b]{0.45\linewidth}
     \includegraphics[width=\linewidth]{Images/RUN_10_figure_2.png}
  \end{subfigure}
  \caption{Comparison of Wind Velocity components from \textit{RUN\_01},
    and \textit{RUN\_10}.}
  \label{fig:RUN10_2}
\end{figure}
\clearpage

%***********************************************************************
% RUN_07
%***********************************************************************
\subsection{RUN\_11\_average\_wind}
\textbf{Configuration:}\\ preload\\
\textit{Binaries/DRWP\_w\_comp.bin}, profile $2$ (\textit{RUN\_02})\\
\textit{Binaries/DRWP\_no\_w\_comp.bin}, profile $1771$ (\textit{RUN\_01})

\textbf{Purpose:}\\ To test the different methods of calling
\textit{compute\_average\_wind(...)}.

The logged output for this test-run includes the same variables as the previous
test-runs, with the addition of the average velocity outputs at the end. It
is important to compare the output of this run with the output from
\textit{RUN\_02} because the computation of average velocity should not
interfere with the data generated to be representative of the current
vehicle, even when the average velocity is being computed from a different
profile. \textbf{The interpolated data values match with those from
\textit{RUN\_02}, as required.}

The average velocity data is shown below:

\includegraphics[width=0.7\linewidth]{Images/data_11.png}

Summary of results:
\begin{enumerate}
\setcounter{enumi}{0}
\item average winds not computed, values are zero-vector
\item Zero-argument method: average winds computed over full domain for
  the current index (0).
\item One-argument method: average winds computed over full domain
  specifically for index 0.
\item Two-argument method: average winds computed over domain [0,1000]
  for index 0.
\item Two-argument method: average winds computed over domain [0,10000]
  for index 0.
\item One-argument method: average winds computed over full domain
  for index 1 while the model continues to use index 0 for scheduled updates.
\item Three-argument method: average winds computed over domain [0,1000]
  for index 1 while the model continues to use index 0 for scheduled updates.
\item Three-argument method: average winds computed over domain [0,10000]
  for index 1 while the model continues to use index 0 for scheduled updates.
\item Two-argument method: average winds computed over domain [0,0.1]
  for index 0. There are no calibrated points in this domain, an error
  is posted and the average velocity is set to the zero vector.
\item Three-argument method: average winds computed over domain [0,1000]
  for index 5. There is no profile loaded to index 5, an error is posted
  and the average velocity is set to the zero vector.
\end{enumerate}


%***********************************************************************
% ERROR cases
%***********************************************************************
\subsection{ERROR\_bad\_binary01\_no\_number\_of\_profiles}
\textbf{Configuration:}\\ preload\\
\textit{Binaries/DRWP\_no\_w\_comp.bin}, profile $1771$ (\textit{RUN\_01})\\
run-time load\\
\textit{Binaries/corrupt\_binaries/no\_number\_of\_profiles.bin}, profile $11$

\textbf{Purpose:}\\ To test the error detection of
a run-time load of an invalid binary file (empty file)
when a valid binary file is already pre-loaded.

\includegraphics[width=\linewidth]{Images/error1.png}


\subsection{ERROR\_bad\_binary03\_no\_number\_of\_altitudes}
\textbf{Configuration:}\\ preload\\
\textit{Binaries/DRWP\_no\_w\_comp.bin}, profile $1771$ (\textit{RUN\_01})\\
run-time load\\
\textit{Binaries/corrupt\_binaries/no\_number\_of\_altitudes.bin}, profile $11$

\textbf{Purpose:}\\ To test the error detection with
a run-time load of an invalid binary file (no number of altitudes entry)
when a valid binary file is already pre-loaded.

\includegraphics[width=\linewidth]{Images/error2.png}



\subsection{ERROR\_bad\_binary05\_short\_altitude\_array}
\textbf{Configuration:}\\ preload\\
\textit{Binaries/DRWP\_no\_w\_comp.bin}, profile $1771$ (\textit{RUN\_01})\\
run-time load\\
\textit{Binaries/corrupt\_binaries/short\_altitude\_array.bin}, profile $11$

\textbf{Purpose:}\\ To test the error detection
a run-time load of an invalid binary file (incorrectly sized altitude array)
when a valid binary file is already pre-loaded.

\includegraphics[width=\linewidth]{Images/error3.png}



\subsection{ERROR\_bad\_binary07\_short\_profile\_array}
\textbf{Configuration:}\\ preload\\
\textit{Binaries/DRWP\_no\_w\_comp.bin}, profile $1771$ (\textit{RUN\_01})\\
run-time load\\
\textit{Binaries/corrupt\_binaries/short\_profile\_array.bin}, profile $11$

\textbf{Purpose:}\\ To test the error detection
a run-time load of an invalid binary file (incorrectly sized profile array)
when a valid binary file is already pre-loaded.

\includegraphics[width=\linewidth]{Images/error4.png}



\subsection{ERROR\_bad\_binary08\_no\_profile\_number}
\textbf{Configuration:}\\ preload\\
\textit{Binaries/DRWP\_no\_w\_comp.bin}, profile $1771$ (\textit{RUN\_01})\\
run-time load\\
\textit{Binaries/corrupt\_binaries/no\_profile\_number.bin}, profile $11$

\textbf{Purpose:}\\ To test the error detection
a run-time load of an invalid binary file (missing profile data)
when a valid binary file is already pre-loaded.

\includegraphics[width=\linewidth]{Images/error5.png}



\subsection{ERROR\_deprecated}
\textbf{Configuration:}\\ preload\\
\textit{Binaries/DRWP\_no\_w\_comp.bin}, profile $1771$ (\textit{RUN\_01})

\textbf{Purpose:}\\ To test the error detection when executing deprecated
methods.

\includegraphics[width=\linewidth]{Images/error6.png}



\subsection{ERROR\_no\_data}
\textbf{Configuration:}\\ Empty

\textbf{Purpose:}\\ To test the error detection when the model is provided no
interpolation data or a configuration to find interpolation data.

\includegraphics[width=\linewidth]{Images/error7.png}



\subsection{FAIL\_01\_repetitive\_preload}
\textbf{Configuration:}\\ preload
\textit{Binaries/DRWP\_no\_w\_comp.bin}, profile $1771$ 11 times

\textbf{Purpose:}\\ To test the error detection when pre-loading more profiles
than slots available.

\includegraphics[width=\linewidth]{Images/fail1.png}



\subsection{FAIL\_06\_spec\_6\_file\_invalid\_profile}
\textbf{Configuration:}\\
run-time load\\
\textit{Binaries/DRWP\_w\_comp.bin}, profile $12$ (invalid).

\textbf{Purpose:}\\ To test the error detection when specifying an invalid
profile number in a good data file.

\includegraphics[width=\linewidth]{Images/fail2.png}


\subsection{FAIL\_06\_spec\_6\_file\_invalid\_vertical}
\textbf{Configuration:}\\
run-time load\\
\textit{Binaries/DRWP\_w\_comp.bin}, profile $2$ (\textit{RUN\_02}), with
specification that profile does not include vertical wind data.

\textbf{Purpose:}\\ To test the error detection when attempting to load a data
profile with an inconsistent setting of \textit{include\_vertical\_component}.

\includegraphics[width=\linewidth]{Images/fail3.png}


\subsection{FAIL\_07\_mid\_run\_switch\_invalid\_index}
\textbf{Configuration:}\\ preload\\
\textit{Binaries/DRWP\_no\_w\_comp.bin}, profile $1771$ (\textit{RUN\_01})

\textbf{Purpose:}\\ To test the error detection when switching to a
pre-loaded profile index higher than the availability of pre-loaded
profiles. In this case, there is 1 pre-loaded profile, and we try to switch
to index 5.

\includegraphics[width=\linewidth]{Images/fail4.png}


\subsection{FAIL\_bad\_binary01\_no\_number\_of\_profiles}
\textbf{Configuration:}\\ preload\\
\textit{Binaries/corrupt\_binaries/no\_number\_of\_profiles.bin}, profile $11$

\textbf{Purpose:}\\ To test the error detection with
a pre-load of an invalid binary file (empty file).

\includegraphics[width=\linewidth]{Images/fail5.png}


\subsection{FAIL\_bad\_binary02\_invalid\_number\_of\_profiles}
\textbf{Configuration:}\\ preload\\
\textit{Binaries/corrupt\_binaries/invalid\_number\_of\_profiles.bin}, profile $11$

\textbf{Purpose:}\\ To test the error detection with
a pre-load of an invalid binary file (invalid number of profiles).

\includegraphics[width=\linewidth]{Images/fail6.png}


\subsection{FAIL\_bad\_binary03\_no\_number\_of\_altitudes}
\textbf{Configuration:}\\ preload\\
\textit{Binaries/corrupt\_binaries/no\_number\_of\_altitudes.bin}, profile $11$

\textbf{Purpose:}\\ To test the error detection with
a pre-load of an invalid binary file (no number of altitudes entry).

\includegraphics[width=\linewidth]{Images/fail7.png}


\subsection{FAIL\_bad\_binary04\_invalid\_number\_of\_altitudes}
\textbf{Configuration:}\\ preload\\
\textit{Binaries/corrupt\_binaries/invalid\_number\_of\_altitudes}, profile $11$

\textbf{Purpose:}\\ To test the error detection with
a pre-load of an invalid binary file (invalid number of altitudes specified).

\includegraphics[width=\linewidth]{Images/fail8.png}

\subsection{FAIL\_bad\_binary05\_short\_altitude\_array}
\textbf{Configuration:}\\ preload\\
\textit{Binaries/corrupt\_binaries/short\_altitude\_array.bin}, profile $11$

\textbf{Purpose:}\\ To test the error detection with
a pre-load of an invalid binary file (incorrectly sized altitude array)

\includegraphics[width=\linewidth]{Images/fail9.png}



\subsection{FAIL\_bad\_binary06\_invalid\_altitude\_array}
\textbf{Configuration:}\\ preload\\
\textit{Binaries/corrupt\_binaries/invalid\_altitude\_array.bin}, profile $11$

\textbf{Purpose:}\\ To test the error detection with
a pre-load of an invalid binary file (non-monotonic data in the altitude array).
This is a compound error: the table-interpolation model detects the problem
with lack of monotonicity and posts an error; this model interprets that
error as being critical and terminates the simulation.

\includegraphics[width=\linewidth]{Images/fail10.png}



\subsection{FAIL\_bad\_binary07\_short\_profile\_array}
\textbf{Configuration:}\\ preload\\
\textit{Binaries/corrupt\_binaries/short\_profile\_array.bin}, profile $11$

\textbf{Purpose:}\\ To test the error detection with
a pre-load of an invalid binary file (incorrectly sized profile array)

\includegraphics[width=\linewidth]{Images/fail11.png}



\subsection{FAIL\_bad\_binary08\_no\_profile\_number}
\textbf{Configuration:}\\ preload\\
\textit{Binaries/corrupt\_binaries/no\_profile\_number.bin}, profile $11$

\textbf{Purpose:}\\ To test the error detection with
a pre-load of an invalid binary file (missing profile data). In this case,
the corrupted binary file is missing the entirety of the profile data,
and the error is detected when the reader fails to read the first part of
that profile, being the profile number.

\includegraphics[width=\linewidth]{Images/fail12.png}



\subsection{FAIL\_bad\_binary09\_invalid\_profile\_number}
\textbf{Configuration:}\\ preload\\
\textit{Binaries/corrupt\_binaries/invalid\_profile\_number.bin}, profile $11$

\textbf{Purpose:}\\ To test the error detection with
a pre-load of an invalid binary file (profile number at the front of the
profile does not match the value in the initial array).

\includegraphics[width=\linewidth]{Images/fail13.png}



\subsection{FAIL\_bad\_binary10\_incorrect\_filelength}
\textbf{Configuration:}\\ preload\\
\textit{Binaries/corrupt\_binaries/incorrect\_filelength.bin}, profile $11$

\textbf{Purpose:}\\ To test the error detection with
a pre-load of an invalid binary file (file length does not match
expectation).

\includegraphics[width=\linewidth]{Images/fail14.png}



\subsection{FAIL\_bad\_filename}
\textbf{Configuration:} Invalid

\textbf{Purpose:}\\ To test the error detection when
an invalid filename is provided for preloading data.

\includegraphics[width=\linewidth]{Images/fail15.png}


\subsection{FAIL\_bad\_filename\_post\_init}
\textbf{Configuration:} Invalid

\textbf{Purpose:}\\ To test the error detection when
an invalid filename is provided for run-time loading of data.
The model starts inactive with no pre-loaded data;
it fails when it is activated with a specified
filename that is not valid.
This scenario results in multiple messages:
\begin{enumerate}
\item At activation, the model searches for pre-loaded
  data matching the specification. On finding none, it posts an inform-level
  message and starts the search for a file with that name.
\item Because this search occurs post-initialization, the
  model also posts an inform-level message about not using recommended
  practice (which is to have all data loaded before initialization).
\item The model is unable to find the specified file; because the model has
  no data already loaded, this results in termination.
\end{enumerate}


\includegraphics[width=\linewidth]{Images/fail16.png}



\subsection{FAIL\_bad\_filename\_with\_good\_preload}
\textbf{Configuration:}\\ preload\\
\textit{Binaries/DRWP\_no\_w\_comp.bin}, profile $1771$ (\textit{RUN\_01})\\
run-time load\\Invalid

\textbf{Purpose:}\\ To test the error detection when
an invalid filename is provided for the initial identification of
the data table to use. This scenario results in multiple messages:
\begin{enumerate}
\item At activation, the model searches for pre-loaded
  data matching the specification. On finding none, it posts an inform-level
  message and starts the search for a file with that name.
\item Because this search occurs post-initialization (activation always
  occurs post-initialization due to the nature of the Subscription model), the
  model also posts an inform-level message about not using recommended
  practice (which is to have all data loaded before initialization).
\item The model is unable to find the specified file; because the model has
  data already loaded, this results in a non-critical error.
\item The failure to read the file means the data fails to load, resulting
  in the profile not being changed, which is a non-critical error.
\item Failure to change the profile during activation means that no profile
  is available for the model to use; this is a critical fault, and the model
  terminates the simulation.
\end{enumerate}

\includegraphics[width=\linewidth]{Images/fail17.png}

\end{document}
